\section{Requerimientos específicos}

\subsection{I. Requisitos funcionales}

Los requisitos funcionales se agrupan por módulo. Cada módulo mantiene su propio archivo
para permitir agregar, retirar o modificar funcionalidades sin alterar los demás.

% Archivos de requisitos funcionales por módulo
\subsubsection{Módulo: [NOMBRE DEL MÓDULO 1]}

\Campo{DESCRIPCIÓN BREVE DEL MÓDULO.}

\begin{longtable}{p{0.22\textwidth} p{0.70\textwidth}}
    \toprule
    \textbf{ID} & \textbf{Requisito funcional}\\
    \midrule
    \endhead
    RF-[MOD]-001 & El sistema deberá \Campo{ACCIÓN O COMPORTAMIENTO ESPERADO.}\\
    RF-[MOD]-002 & El sistema permitirá \Campo{ACCIÓN O COMPORTAMIENTO ESPERADO.}\\
    RF-[MOD]-003 & El sistema deberá \Campo{ACCIÓN O COMPORTAMIENTO ESPERADO.}\\
    \bottomrule
\end{longtable}

\subsubsection{Módulo: [NOMBRE DEL MÓDULO 2]}

\Campo{DESCRIPCIÓN BREVE DEL MÓDULO.}

\begin{longtable}{p{0.22\textwidth} p{0.70\textwidth}}
    \toprule
    \textbf{ID} & \textbf{Requisito funcional}\\
    \midrule
    \endhead
    RF-[MOD]-001 & El sistema deberá \Campo{ACCIÓN O COMPORTAMIENTO ESPERADO.}\\
    RF-[MOD]-002 & El sistema permitirá \Campo{ACCIÓN O COMPORTAMIENTO ESPERADO.}\\
    \bottomrule
\end{longtable}

% Para agregar un módulo:
% 1. Copiar este archivo y asignarle un nombre nuevo.
% 2. Cambiar el prefijo [MOD] por el identificador del módulo.
% 3. Agregar \input{modulos/nombre-del-archivo} en la sección 3.


\subsection{II. Requisitos de interfaces externas}

\subsubsection{Interfaces de usuario}
\begin{longtable}{p{0.22\textwidth} p{0.70\textwidth}}
    \toprule
    \textbf{ID} & \textbf{Requisito propuesto}\\
    \midrule
    \endhead
    RI-UI-001 & \Campo{DESCRIPCIÓN DEL REQUISITO.}\\
    RI-UI-002 & \Campo{DESCRIPCIÓN DEL REQUISITO.}\\
    \bottomrule
\end{longtable}

\subsubsection{Interfaces de software}
\begin{longtable}{p{0.22\textwidth} p{0.70\textwidth}}
    \toprule
    \textbf{ID} & \textbf{Requisito propuesto}\\
    \midrule
    \endhead
    RI-SW-001 & \Campo{DESCRIPCIÓN DEL REQUISITO.}\\
    RI-SW-002 & \Campo{DESCRIPCIÓN DEL REQUISITO.}\\
    \bottomrule
\end{longtable}

\subsubsection{Interfaces de hardware}
\begin{longtable}{p{0.22\textwidth} p{0.70\textwidth}}
    \toprule
    \textbf{ID} & \textbf{Requisito propuesto}\\
    \midrule
    \endhead
    RI-HW-001 & \Campo{DESCRIPCIÓN DEL REQUISITO.}\\
    RI-HW-002 & \Campo{DESCRIPCIÓN DEL REQUISITO.}\\
    \bottomrule
\end{longtable}

\subsection{III. Requisitos no funcionales}
\begin{longtable}{p{0.20\textwidth} p{0.20\textwidth} p{0.52\textwidth}}
    \toprule
    \textbf{ID} & \textbf{Categoría} & \textbf{Requisito propuesto}\\
    \midrule
    \endhead
    RNF-COM-001 & Compatibilidad & \Campo{DESCRIPCIÓN DEL REQUISITO.}\\
    RNF-REN-001 & Rendimiento & \Campo{DESCRIPCIÓN DEL REQUISITO.}\\
    RNF-SEG-001 & Seguridad & \Campo{DESCRIPCIÓN DEL REQUISITO.}\\
    RNF-USA-001 & Usabilidad & \Campo{DESCRIPCIÓN DEL REQUISITO.}\\
    RNF-MAN-001 & Mantenibilidad & \Campo{DESCRIPCIÓN DEL REQUISITO.}\\
    \bottomrule
\end{longtable}
